\documentclass[10pt]{article}       % set main text size
\usepackage[letterpaper,                % set paper size to letterpaper. change to a4paper for resumes outside of North America
top=0.5in,                          % specify top page margin
bottom=0.5in,                       % specify bottom page margin
left=0.5in,                         % specify left page margin
right=0.5in]{geometry}              % specify right page margin
                       
\usepackage{XCharter}               % set font. comment this line out if you want to use the default LaTeX font Computer Modern
\usepackage[T1]{fontenc}            % output encoding
\usepackage[utf8]{inputenc}         % input encoding
\usepackage{enumitem}               % enable lists for bullet points: itemize and \item
\usepackage[hidelinks]{hyperref}    % format hyperlinks
\usepackage{titlesec}               % enable section title customization
\raggedright                        % disable text justification
\pagestyle{empty}                   % disable page numbering

% ensure PDF output will be all-Unicode and machine-readable
\input{glyphtounicode}
\pdfgentounicode=1

% format section headings: bolding, size, white space above and below
\titleformat{\section}{\bfseries\large}{}{0pt}{}[\vspace{1pt}\titlerule\vspace{-6.5pt}]

% format bullet points: size, white space above and below, white space between bullets
\renewcommand\labelitemi{$\vcenter{\hbox{\small$\bullet$}}$}
\setlist[itemize]{itemsep=-2pt, leftmargin=12pt, topsep=6pt} %%% Test various topsep values to fix vertical spacing errors

% resume starts here
\begin{document}

% name
\centerline{\Huge Chad Gauthier}

% \vspace{5pt}

% contact information
\centerline{\href{mailto:chad.gauthier32@gmail.com}{chad.gauthier32@gmail.com} | \href{https://www.github.com/cgoat24}{github.com/cgoat24}}

\vspace{-18.5pt}

% education section
\section*{Éducation}
\textbf{Université de Sherbrooke} -- Baccalauréat en informatique \hfill 2025 \\
\textbf{Cégep de Trois-Rivières} -- Diplôme d'étude collégiale en informatique \hfill 2022

\vspace{-6.5pt}

% experience section
\section*{Expérience}
\textbf{Programmeur-analyste stagiaire} {Nmédia} \hfill Mai 2024 -- Août 2024 \\
\vspace{-9pt}
\begin{itemize}
  \item Créé une plateforme d'analyse de données pour évaluer les performances des solutions et soutenir les décisions basées sur les données.
  \item Refactorisé divers modules pour améliorer la stabilité, la maintenabilité et les performances de la base de code.
  \item Collaboré avec une équipe agile pour apporter des améliorations itératives, en participant activement à la planification des sprints, aux réunions quotidiennes et aux rétrospectives.
\end{itemize}

\textbf{Programmeur outil stagiaire} {Ubisoft Québec} \hfill Septembre 2023 -- Décembre 2023 \\
\vspace{-9pt}
\begin{itemize}
	\item Conçu des outils internes pour automatiser les tâches, augmenter l'efficacité lors du développement de jeux.
	\item Implémenté des fonctionnalités dans une application web en Go, améliorant l'expérience utilisateur.
	\item Écrit des tests de bout en bout avec Cypress pour valider les flux d'utilisateurs critiques.
\end{itemize}

\textbf{Programmeur-analyste stagiaire} {Sherweb} \hfill Janvier 2023 -- Avril 2023 \\
\vspace{-9pt}
\begin{itemize}
	\item Implémentation de nouvelles fonctionnalités dans une application de type microservices déployée sur Kubernetes, permettant une orchestration de services évolutive et résiliente.
	\item Réalisé un outil interne pour rationaliser les flux de travail et améliorer l'expérience des développeurs, réduisant ainsi les frictions dans les tâches d'ingénierie quotidiennes.
	\item Collaboré avec une équipe agile pour apporter des améliorations itératives, en participant activement à la planification des sprints, aux réunions quotidiennes et aux rétrospectives
\end{itemize}

\textbf{Programmeur-analyste} {Groupe Shift} \hfill Mars 2022 -- Septembre 2022 \\
\vspace{-9pt}
\begin{itemize}
	\item Contribué à plus de 5 projets simultanément, en mettant en œuvre des solutions selon les exigences du client.
	\item Participé à un comité de recherche axé sur le développement SaaS, en fournissant des informations et des recommandations pour soutenir les décisions stratégiques en matière de produits.
	\item Collaboré avec des équipes interfonctionnelles pour recueillir les spécifications, documenter le comportement du système et valider les besoins des clients.
\end{itemize}

\vspace{-18.5pt}

% projects section
\section*{Projets}
\textbf{FlexCalc} \hfill \href{https://github.com/cgoat24/flex-calc}{github.com/cgoat24/flex-calc} \\
\vspace{-9pt}
\begin{itemize}
  \item Développé une application mobile destinée aux passionnés de musculation à l'aide de Svelte, Rust et Tauri.
  \item Implémenté une logique de calcul prenant en charge un large éventail d'opérations mathématiques.
  \item Assuré une interface conviviale et une intégration transparente de nouvelles fonctionnalités.
\end{itemize}

\textbf{TimeTunes} \hfill \href{https://github.com/cgoat24/timetunes}{github.com/cgoat24/timetunes} \\
\vspace{-9pt}
\begin{itemize}
  \item Codéveloppé un jeu mobile multijoueur permettant aux joueurs de deviner les dates de sortie des chansons en les plaçant sur une chronologie.
  \item Intégré le SDK Spotify pour diffuser de la musique directement dans l'application, permettant aux utilisateurs de jouer avec n'importe quelle liste de lecture publique.
  \item Implémenté une fonctionnalité multijoueur en temps réel, permettant aux utilisateurs de jouer à distance avec des amis, chacun sur son propre appareil.
\end{itemize}

\textbf{La Beuverie} \hfill \href{https://la-beuverie.vercel.app}{la-beuverie.vercel.app} \\
\vspace{-9pt}
\begin{itemize}
  \item Conçu et développé une application Web complète pour partager des cocktails, utilisant Svelte et Typescript.
  \item Configuration de Docker pour le développement local, permettant des environnements reproductibles.
  \item Déploiement sur Vercel, garantissant une livraison évolutive et performante de l'app aux utilisateurs.
\end{itemize}

\vspace{-18.5pt}

% skills section
\section*{Compétences}
\textbf{Langages:} C\#, Javscript, Typescript, Rust, Python, Go, PHP, Java, C/C++ \\
\textbf{Librairies:} Vue.js, Nuxt, Svelte, Sveltekit, React, Astro, TailwindCSS, .NET, Laravel, Express.js, Hono, Prisma \\
\textbf{Bases de données:} SQL Server, PostgreSQL, MongoDb, MySQL \\
\textbf{Outils:} Docker, Git, Github, Gitlab, Vercel, AWS, Linux

\end{document}